\section{Исходный код программы}
\label{app:source}
Программа находится в файле \texttt{code/search.py}. Для запуска
из корня проекта используется команда \verb|python3 code/search.py|.
Разделитель столбцов в выводе --- точка с запятой.

\captionof{listing}{Функции поиска}
\label{code:search}
\inputminted[firstline=1,lastline=24]{python}{code/search.py}

\clearpage
\captionof{listing}{Запуск эксперимента}
\label{code:experiment}
\inputminted[firstline=27,firstnumber=27]{python}{code/search.py}

Чтобы пересоздать изображения, можно выполнить
\verb|python3 code/plots.py|. Для этого нужен пакет \texttt{matplotlib}.
Для сборки документа он не требуется: готовые изображения уже находятся
в папке \texttt{images}. Исходный код построения приведён
в листинге~\ref{code:plots}.
\clearpage
\captionof{listing}{Построение графиков}
\label{code:plots}
\inputminted{python}{code/plots.py}
